\documentclass[aspectratio=169]{beamer}
\usetheme{default}
\setbeamertemplate{navigation symbols}{}
\usepackage{amsmath,amssymb}
\title{Kodaira--Spencer and the tangent space of $\mathrm{Def}(X)$}
\subtitle{Huybrechts, \S1}
\date{Reading seminar}
\begin{document}
\begin{frame}\titlepage\end{frame}
\begin{frame}{1. Definition of a deformation}
Following Huybrechts, a deformation of a compact complex
manifold $X$ is a proper smooth holomorphic family
\[
\pi:\mathcal X\longrightarrow B,\qquad
0\in B,\quad \pi^{-1}(0)\simeq X,
\]
considered near $0$. The fibres $X_t$ vary the complex structure.

A tangent vector $v\in T_0B$ is an infinitesimal direction.
The Kodaira--Spencer map turns it into a class on $X$.
\end{frame}

\begin{frame}{2. The relative tangent sequence}
Because $\pi$ is a submersion, its differential is surjective.
The kernel consists of vectors tangent to the fibres:
\[
0\longrightarrow T_{\mathcal X/B}
\longrightarrow T_{\mathcal X}
\xrightarrow{d\pi}\pi^*T_B\longrightarrow0.
\]
Restrict to the central fibre. Since
$T_{\mathcal X/B}|_X=T_X$ and
$\pi^*T_B|_X=T_0B\otimes\mathcal O_X$, we obtain
\[
\boxed{0\longrightarrow T_X\longrightarrow
T_{\mathcal X}|_X\longrightarrow
T_0B\otimes\mathcal O_X\longrightarrow0.}
\]
\end{frame}

\begin{frame}{3. The connecting homomorphism trick}
Apply the long exact sequence of sheaf cohomology:
\[
H^0(X,T_{\mathcal X}|_X)\longrightarrow
\underbrace{H^0(X,T_0B\otimes\mathcal O_X)}_{=\,T_0B}
\xrightarrow{\ \delta\ }H^1(X,T_X)\longrightarrow\cdots
\]
Here $H^0(X,\mathcal O_X)=\mathbb C$ for compact connected $X$.
The connecting map is precisely Kodaira--Spencer:
\[
\boxed{\rho=\delta:T_0B\longrightarrow H^1(X,T_X).}
\]
A local lift of $v$ to $T_{\mathcal X}|_X$ need not glue
globally. Differences of local lifts are tangent to $X$
and define a Cech $1$-cocycle representing $\rho(v)$.
\end{frame}
\begin{frame}{4. The transition-function picture (HP)}
Choose local trivialisations with changes of coordinates
$(x,t)\mapsto(\psi_{ij}(x,t),t)$.
For $v\in T_0B$, differentiating at $t=0$ gives
\[
\theta_{ij}=\left.\frac d{d\varepsilon}\right|_0
\psi_{ij}(x,\varepsilon v).
\]
These vector fields give a Cech $1$-cocycle;
changing trivialisations changes it by a coboundary.
Thus $\rho(v)=[\theta_{ij}]\in H^1(X,T_X)$.
\smallskip

{\scriptsize From HP \texttt{deformations.tex}, Kodaira--Spencer map;
\texttt{complex-geometry.tex}, relative tangent sequence.}
\end{frame}
\begin{frame}{5. The IHS computation}
Dolbeault gives $H^1(X,T_X)\simeq
H_{\bar\partial}^{0,1}(X,T^{1,0}X)$.
For an IHS manifold, contraction with the nondegenerate
holomorphic symplectic form $\sigma$ is an isomorphism
\[
T_X\xrightarrow{\ \iota_{(-)}\sigma\ }\Omega_X^1.
\]
Hence, by Dolbeault and the Kahler Hodge decomposition,
\[
\boxed{T_0\operatorname{Def}(X)
\simeq H^1(X,T_X)\simeq H^1(X,\Omega_X^1)
\simeq H^{1,1}(X).}
\]
The Bogomolov--Tian--Todorov theorem ensures that
$\operatorname{Def}(X)$ is smooth.
\end{frame}
\end{document}
